\section{This Water, This Bread}

\subsection{The particularity did not retire at the Ascension}

A tempting compromise would concede everything so far and then quietly re-abstract the result: granted, God acted particularly \emph{then}; but now the Incarnation survives as a truth to be believed, an example to be admired---portable, general, safely conceptual. That is, in effect, Kant's proposal with the dates filled in: the husk has done its work and may now split. The Church's sacramental teaching exists to block precisely that retreat. Leo the Great, preaching on the Ascension, states the principle in a single clause: ``that which till then was visible of our Redeemer was changed into a sacramental presence.''\endnote{Leo the Great, \emph{Sermo} 74 (De ascensione Domini II), 2, in the public-domain Nicene and Post-Nicene Fathers translation (Charles Lett Feltoe), checked at the New Advent transcription, 24 July 2026.} The visible body departs; the mode of presence changes; the particularity does not. What replaces the carpenter's hands is not an idea but a set of instituted, perceptible acts---washing, anointing, the giving of bread and the cup---each as concrete, as datable, as located as Nazareth.

The Catechism explains why this is not a lapse back into magic but the continuation of the same divine style. The sacraments are ``perceptible signs (words and actions) accessible to our human nature,'' by which the risen Christ, seated at the Father's right hand, now acts to communicate the grace they signify. And they are possible because the Paschal event is unlike every other event: all others ``happen once, and then they pass away, swallowed up in the past,'' but the Paschal mystery of Christ ``cannot remain only in the past,'' since in it death itself was destroyed and everything Christ did and suffered participates in divine eternity.\endnote{\emph{Catechism of the Catholic Church} 1084--1085, official English text at the Vatican website, checked 24 July 2026. This is a claim of sacramental making-present, not of repetition: the historical event is not reenacted but its saving power made available. The article does not expand here into a treatise on sacramental validity, matter and form, or the theology of Eucharistic presence, all of which have their own governing sources.} Lessing's ditch was dug on the assumption that a past event can reach the present only by report, dwindling as testimony dwindles. The sacramental economy is the claim that this one event reaches the present by \emph{presence}---that the ditch is crossed, not by our leap backward, but by the event's own refusal to recede.

\subsection{Why signs suit us}

Aquinas's defense of sensible signs completes the pattern this article has been tracing since the discussion of abstraction, and it is worth taking in its full structure rather than in a single line.

He begins from what a sign is and whom signs are for: ``Signs are given to men, to whom it is proper to discover the unknown by means of the known.''\endnote{Aquinas, \emph{Summa theologiae} III, q.~60, a.~2, corpus, English Dominican translation, checked 25 July 2026. The article is defining a sacrament as ``the sign of a holy thing so far as it makes men holy,'' and the clause quoted supplies the anthropological premise on which the rest of the question turns.} Then the governing principle: divine wisdom ``provides for each thing according to its mode,'' and it is the human mode---the point was established about knowers, and returns here about worshipers---``to acquire knowledge of the intelligible from the sensible.'' Hence sacraments are sensible things; and not sensible things in general, but ``those things which are determined by Divine institution,'' since ``it is not for anyone to determine that which is in the power of another,'' and sanctification is in God's power, not ours.\endnote{Aquinas, \emph{Summa theologiae} III, q.~60, aa.~4--5, corpus, English Dominican translation, checked 24 July 2026 and re-read 25 July 2026. Article~4's \emph{sed contra} appeals to Wisdom 8:1 and to Dionysius; article~5's argument is that determining the means of another's act belongs to that other. The determinacy point matters to this article's argument and is why it is quoted: an economy of instituted signs is more particular, not less, than an economy of natural symbols.}

That last point deserves more weight than it usually gets. The signs are not merely sensible; they are \emph{determined}. Water, not any cleansing liquid; bread and wine, not any meal. A religion of natural symbolism would let each culture select its own emblems, which is precisely what the pluralist reading of religion expects and what Celsus's superintending spirits would have arranged. The instituted economy does the opposite: it fixes the particulars, and it fixes them by an act of choice on the giver's side. This is the same grammar as election, transposed from peoples to elements.

Aquinas then adds an argument that ties the sacraments directly to the Incarnation, and it is the keystone of this whole section. Why must words be added to the sensible element? Because the sacraments must have ``a certain conformity'' to the cause of sanctification, ``which is the Word incarnate \ldots\ in that the word is joined to the sensible sign, just as in the mystery of Incarnation the Word of God is united to sensible flesh.''\endnote{Aquinas, \emph{Summa theologiae} III, q.~60, a.~6, corpus, English Dominican translation, checked 25 July 2026. Aquinas gives three reasons for the addition of words; only the first, the conformity to the Incarnation, is used here. The second---that the remedy is fitted to a creature composed of soul and body, ``since it touches the body through the sensible element, and the soul through faith in the words''---bears on the following paragraph.} The structure of a sacrament, on this account, is not merely convenient; it is a deliberate reprise. The pattern of word-joined-to-flesh, having happened once at Nazareth, is made repeatable at every font. If one finds the Incarnation's particularity offensive, one will find the sacraments offensive for exactly the same reason, and consistently so. The tradition, for its part, treats the resemblance as an argument in favor.

The necessity of such signs, Aquinas argues in the following question, has three roots, and the second is the one this article most needs. The first is the condition of human nature, which ``has to be led by things corporeal and sensible to things spiritual and intelligible.'' The second is the state of fallen man, ``who in sinning subjected himself by his affections to corporeal things''---and here the principle is medical: ``the healing remedy should be given to a man so as to reach the part affected by disease.''\endnote{Aquinas, \emph{Summa theologiae} III, q.~61, a.~1, corpus, English Dominican translation, checked 25 July 2026. The third reason---that bodily exercise in the sacraments displaces harmful and superstitious practice---is Aquinas's own and is not relied on here. Aquinas's reply to the second objection states the boundary this article observes: ``God's grace is a sufficient cause of man's salvation. But God gives grace to man in a way which is suitable to him.''} That is Athanasius's argument about the small instrument, restated for the age after the Ascension. The repair goes where the wound is; the wound was in the flesh; therefore the repair is applied through flesh. A purely spiritual religion---all interiority, no matter---would be fitted for angels, and its persistent attraction for us is one more instance of the flight into abstraction this article has been resisting throughout: the preference for a God who makes no landings. The actual economy is humbler and more invasive. Grace comes to a creature of flesh through flesh: through water that is wet, bread that is bread, words said aloud by one particular minister to one particular person with a name.

\subsection{Signs of the covenant}

The Catechism sets out the same economy historically, and the sequence it describes is the article's argument in liturgical form. Sacramental signs, it says, have a meaning ``rooted in the work of creation and in human culture, specified by the events of the Old Covenant and fully revealed in the person and work of Christ.'' Three strata, each more particular than the last.

The first stratum is universal and natural. ``As a being at once body and spirit, man expresses and perceives spiritual realities through physical signs and symbols''; and God ``speaks to man through the visible creation,'' in which ``light and darkness, wind and fire, water and earth, the tree and its fruit speak of God and symbolize both his greatness and his nearness.'' The Catechism then says something the objector should notice, since it concedes a great deal of what the pluralist wants: ``The great religions of mankind witness, often impressively, to this cosmic and symbolic meaning of religious rites.''\endnote{\emph{Catechism of the Catholic Church} 1145--1149, official English text, checked at the Vatican website, 25 July 2026. The concession at 1149 is the Catechism's own and is quoted here because a fair statement of the Catholic position must include it; the same paragraph adds that the Church's liturgy ``presupposes, integrates and sanctifies elements from creation and human culture, conferring on them the dignity of signs of grace.'' The relation of the Church to other religions is treated at its own place below and under its own governing documents.}

The second stratum is where particularity enters, and the Catechism names it exactly: ``Signs of the covenant.'' The chosen people ``received from God distinctive signs and symbols that marked its liturgical life. These are no longer solely celebrations of cosmic cycles and social gestures, but signs of the covenant, symbols of God's mighty deeds for his people''---circumcision, the anointing of kings and priests, the laying on of hands, sacrifices, ``and above all the Passover.''\endnote{\emph{Catechism of the Catholic Church} 1150, official English text, checked 25 July 2026. The typological reading that follows in the same paragraph---that the Church sees in these signs ``a prefiguring of the sacraments of the New Covenant''---is the Catechism's own and is reported, not independently argued.} That is the decisive transition, and it is the same transition Deuteronomy made. A cosmic symbol says something true about God in general; a sign of the covenant says something that happened. Water can mean cleansing anywhere; a Passover means a night.

The third stratum narrows again to a person: signs ``taken up by Christ,'' who gives new meaning to the deeds and signs of the Old Covenant, ``for he himself is the meaning of all these signs.'' And then the sacramental signs proper, through which, since Pentecost, the Spirit ``carries on the work of sanctification''---signs which ``do not abolish but purify and integrate all the richness of the signs and symbols of the cosmos and of social life.''\endnote{\emph{Catechism of the Catholic Church} 1151--1152, official English text, checked 25 July 2026.} Cosmos, covenant, Christ, Church: at each stage the sign becomes more particular and the reach becomes wider, which is the pattern this article has been documenting since Abraham.

\subsection{This bread}

The Eucharist is where the pattern reaches its limit, and where the objection, if it is going to break, breaks.

Aquinas's account of why Christ's body is present in the sacrament turns on the same premise as the Incarnation. It ``belongs to Christ's love, out of which for our salvation He assumed a true body of our nature''; and because friendship wants presence, ``meanwhile in our pilgrimage He does not deprive us of His bodily presence; but unites us with Himself in this sacrament through the truth of His body and blood.''\endnote{Aquinas, \emph{Summa theologiae} III, q.~75, a.~1, corpus, English Dominican translation, checked 25 July 2026. Aquinas gives three reasons; the article quotes the second. The first (the perfection of the New Law over against the figures of the Old) and the third (the perfection of faith, ``since faith is of things unseen'') are noted but not relied on. This article states the doctrine of the real presence as the Church's teaching and does not undertake its exposition or defense; sacramental theology proper, including transubstantiation and the conditions of validity, is outside its scope, as its terminal appendix records.} The Eucharist, on this account, is not a departure from the logic of the Incarnation but its continuation under changed conditions: the same love that would not save humanity at a distance will not accompany it at a distance either.

And then a corollary that inverts the objection's arithmetic completely. Christ is present in the sacrament after the manner of substance, not of dimensive quantity; and since ``the whole nature of a substance is under every part of the dimensions under which it is contained''---as the whole nature of bread is under every part of the bread---``it is manifest that the entire Christ is under every part of the species of the bread, even while the host remains entire, and not merely when it is broken.''\endnote{Aquinas, \emph{Summa theologiae} III, q.~76, a.~3, corpus, English Dominican translation, checked 25 July 2026. Aquinas rejects, in the same corpus, the mirror analogy sometimes offered for this, on the ground that the comparison is imperfect. The metaphysics of substance and quantity presupposed here belongs to the sacramental theology this article does not undertake; the point taken is only the formal one stated in the following sentence.} Celsus wanted many bodies so that the gift would reach many places. The doctrine he was mocking says that the whole is in each part---which is precisely what a distributed rationing of the divine could never manage, since rations divide. Whatever one makes of the metaphysics, the shape of the claim is the exact contrary of the one the objection assumed it was attacking. Particularity here is not the fragment left over after the whole has been divided. It is the mode in which the whole is had.

\subsection{The Church's own particularity}

The same grammar governs the community that results, and the New Testament states it in a formula that reads like a deliberate provocation to universalism: ``One body and one Spirit; as you are called in one hope of your calling. One Lord, one faith, one baptism. One God and Father of all, who is above all, and through all, and in us all.''\endnote{Ephesians 4:4--6, Douay--Rheims, checked 25 July 2026. The article cites the passage for the juxtaposition it makes---a series of determinate singulars terminating in the Father of all---and does not enter the ecclesiological questions about the unity of the Church that the passage has been used to argue, which are governed by their own sources.} Count the singulars: one body, one Spirit, one hope, one Lord, one faith, one baptism. Then the sentence's destination: ``one God and Father of all.'' The universality is stated at the end, and everything leading to it is a determinate singular. The passage does not treat this as a tension to be managed; it treats the series as the way the universal claim is made.

Baptism in particular is the hardest case for the objection and the clearest instance of the pattern. It is not a disposition, an insight, or an allegiance held inwardly. It is an act, performed at a place, by someone, on someone, with water, once. There is no general baptism; there are only baptisms. And the offer is universal in exactly the way the earlier sections described: not by being abstract, but by being repeatable.

For the recipient, this particularity is not a limit on assurance but its instrument. An abstract salvation must be verified abstractly---by introspection, by the quality of one's feelings, a notoriously bottomless audit. A sacrament happens or it does not. It happened to \emph{you}, on a date, at a font one can drive to. The definite articles that scandalized Nazareth turn out, at the receiving end, to be the whole comfort: not humanity in general invited to a possible mercy, but this person, washed on this morning, in this parish, by name.

The universal reach that the objector feared was being lost is, once again, achieved rather than abandoned. Precisely because the sacramental acts are small, repeatable, and material, they can be performed anywhere humans are---in cathedrals and prison cells, over any water, in every language---until the particular address has been multiplied toward every particular person. A universal offer that remained abstract would reach no one; the concrete sign, endlessly reiterated, is how the once-for-all event goes out ``even to the farthest part of the earth.''
